\section{From Jacobi identities to a relative physical law}
\label{sec:v9-operators}
The proof of the relative law uses three different types of information.
The first is finite-dimensional algebra, the second is localization of
nonlinear perturbations, and the third is stability of a physical
sublevel integral. We separate these statements before deriving a
nonlinear compatibility law from their combination.

\subsection{The classical operator input}
For a periodic discrete Lagrangian orbit, the discrete Hill formula
of Bolotin and Treschev~\cite[Theorem~2.1]{Hill} relates
$\det(P-I)$ to the periodic action Hessian divided by the product of
mixed edge determinants. The Hessian in our endpoint problem instead
has Dirichlet boundary conditions. Its relevant scalar identity is the
following elementary Schur-complement formula. We include its proof
to specify the boundary convention, not as a new determinant identity.

\begin{proposition}[Dirichlet cofactor and normalization]
\label{prop:v9-cofactor}
Let $L(x_0,\ldots,x_j)=\sum_{i=0}^{j-1}\ell_i(x_i,x_{i+1})$
have a smooth stationary interior at fixed endpoints $x_0=u,x_j=v$.
Suppose its interior Hessian $H$ is positive definite and tridiagonal,
with mixed edge derivatives $\ell_{i,12}=-J_i<0$.
For the stationary value $W(u,v)$,
\begin{equation}\label{eq:v9-cofactor}
 -W_{uv}=\frac{\prod_{i=0}^{j-1}J_i}{\det H}.
\end{equation}
For a reference interior Hessian $H^0$ and reference edges $J_i^0$,
put $G^0=(H^0)^{-1}$ and $\Delta H=H-H^0$. Then
\begin{equation}\label{eq:v9-normalized-cofactor}
 \frac{-W_{uv}}{(-W_{uv})^0}
  =\frac{\prod_i(J_i/J_i^0)}{\det(I+G^0\Delta H)}.
\end{equation}
If $H=D(I-T)$, $D$ is diagonal with $D_{ii}\ge d_0>0$,
$T$ has only its two adjacent diagonals nonzero, and
$\|T\|_{\ell^\infty\to\ell^\infty}\le q<1$, then
\begin{equation}\label{eq:v9-classical-green}
 |(H^{-1})_{ik}|\le \frac{q^{|i-k|}}{d_0(1-q)}.
\end{equation}
These estimates have constants independent of the number of sites.
\end{proposition}
\begin{proof}
For $j\ge2$, differentiating the stationary equations eliminates the
interior and gives
$W_{uv}=-J_0J_{j-1}(H^{-1})_{1,j-1}$.
The corner cofactor of a tridiagonal matrix is
$J_1\cdots J_{j-2}$, so Cramer's rule proves
\eqref{eq:v9-cofactor}. At $j=1$ the interior determinant is the
empty determinant, equal to one, and the same assertion is immediate.
Taking the ratio at the reference point proves
\eqref{eq:v9-normalized-cofactor}. Finally,
$H^{-1}=\sum_{n\ge0}T^nD^{-1}$. The $(i,k)$ entry of $T^n$
vanishes for $n<|i-k|$, while every entry has absolute value at most
$q^n$. Summing the remaining series proves
\eqref{eq:v9-classical-green}.
\end{proof}

The separated inverse structure is classical: Meurant's
Theorem~2.3~\cite{Meurant} gives inverse entries by two scalar
recurrences, and his Theorem~2.4 describes their monotone decrease
under strict diagonal dominance. Nabben~\cite{Nabben} studies
quantitative inverse decay for nonsymmetric tridiagonal and band
matrices. Our use of a weighted Green operator belongs to this
operator framework. In particular neither the corner cofactor nor
exponential inverse decay is claimed here as a new mechanism.

What must additionally be proved for the present probability problem
is summability of the \emph{nonlinear perturbation evaluated on the
entire stationary bridge}. A bound on each matrix entry, without its
location, would not give the required determinant limit. The relevant
continuity statement is recorded next.

\begin{proposition}[A dimension-independent relative determinant bound]
\label{prop:v9-trace-transport}
For trace-class operators $T,T'$ on a common Hilbert space with
$\|T\|,\|T'\|\le q<1$, define
\[
 \mathcal L(T)=\sum_{n\ge1}\frac{(-1)^{n+1}}n\tr(T^n).
\]
Then
\begin{equation}\label{eq:v9-trace-transport}
 |\mathcal L(T)-\mathcal L(T')|
                  \le (1-q)^{-1}\|T-T'\|_1.
\end{equation}
Suppose $T,T'$ depend smoothly on finitely many parameters, their
trace-class derivatives through order $k$ are uniformly bounded,
and the corresponding derivatives of $T-T'$ have trace norm at most
$\epsilon$. Then
\[
       \|\mathcal L(T)-\mathcal L(T')\|_{C^k}\le C_k\epsilon.
\]
The constants depend on $q$, the parameter dimension and the stated
bounds, not on the dimension of a finite-dimensional truncation.
\end{proposition}
\begin{proof}
The algebraic telescoping formula for $T^n-(T')^n$ and the trace ideal
inequality give
$|\tr(T^n-(T')^n)|\le n q^{n-1}\|T-T'\|_1$.
Division by $n$ and summation prove the first assertion.
After $k$ parameter derivatives, each product has at most $k$
differentiated factors. In the difference of the two differentiated
products, telescope the factors once more and place the difference
factor in trace norm. Every other differentiated factor has a bounded
operator norm. All but at most $k+1$ of the factors still have norm
at most $q$. There are at most $C_k n^{k+1}$ terms. The finitely many
products with $n\le k+1$ are bounded separately; the other terms are
bounded by $C_k\epsilon n^{k+1}q^{n-k-1}$. Their sum, including
$1/n$, converges. The same estimates justify differentiating the
original series by uniform convergence. No smallness of a
differentiated operator has been used.
\end{proof}

In the billiard application, Lemma~\ref{lem:g-relative} gives
$|x_i|\le C(|u|\rho^i+|v|\rho^{j-i})$, with the corresponding
localized derivative bounds. Since a Hessian entry depends on only
neighboring coordinates, its perturbation has a uniformly summable
entrywise absolute value. Matrix units have trace norm one; hence
\[
 \|\Delta H_j\|_1\le C(|u|+|v|),\qquad
 \|\partial^\alpha\Delta H_j\|_1\le C_\alpha.
\]
The two-block construction in Theorem~\ref{thm:v4-factorization}
then compares the normalized finite determinant with the direct sum
of the two half-line determinants. Its trace-norm error tends to zero
exponentially. Proposition~\ref{prop:v9-trace-transport} converts that
specific geometric estimate into convergence of the logarithmic
amplitude. Inverse decay supplies access to the blocks; the localized
perturbation estimate supplies their summable determinant error.

\subsection{The passage through physical integration}
A second, logically separate requirement is uniformity as the offset
vanishes. The following statement isolates this passage. All integrals
are over a fixed neighborhood containing the indicated sublevel sets.
The parameter set is compact and finite dimensional.

\begin{proposition}[Stability of normalized sublevel integration]
\label{prop:v9-morse-transport}
Let $E,\widetilde E$ be smooth functions of $y\in\R^2$ and of
parameters, with value and $y$-gradient zero at $y=0$. Assume their
Hessians at zero are uniformly positive and their higher derivatives
are uniformly bounded on a common neighborhood. Let $a,\widetilde a$
be smooth amplitudes with uniform derivative bounds. There is a common
$d_*>0$ such that
\[
 \mathcal J(E,a)(d)=d^{-2}\int(d-E(y))_+a(y)\,\dd y
\]
extends smoothly from $d>0$ to $[0,d_*]$. For each fixed $k$, bounds
through a sufficiently high finite order $r(k)$ give
\begin{equation}\label{eq:v9-morse-transport}
 \|\mathcal J(E,a)-\mathcal J(\widetilde E,\widetilde a)\|_{C^k}
 \le C_k\bigl(\|E-\widetilde E\|_{C^{r(k)}}
                    +\|a-\widetilde a\|_{C^{r(k)}}\bigr).
\end{equation}
The norms on the left include mixed parameter and offset derivatives.
\end{proposition}
\begin{proof}
Taylor's integral formula gives $E(y)=y^tM_E(y)y/2$, with
$M_E(y)=2\int_0^1(1-s)D^2E(sy)\,\dd s$.
After decreasing the common neighborhood, the maps
$y\mapsto M_E(y)^{1/2}y$ and their counterparts for $\widetilde E$
have uniformly invertible differentials. Their inverses $\Phi_E$
are obtained by the contraction used in Lemma~\ref{lem:g-radial}.
Subtracting the inverse equations gives their Lipschitz dependence
in $C^0$ on $E$; differentiating and isolating the highest derivative
on the same uniformly invertible differential proves the corresponding
finite-order estimates by induction. Smoothness of the positive matrix
square root is used only on a compact positive-definite set.

Set $A_E(w)=a(\Phi_E(w))|\det D\Phi_E(w)|$. The change of variables
$w=\sqrt{2d}\,z$ gives the fixed-domain formula
\begin{equation}\label{eq:v9-fixed-disk}
 \mathcal J(E,a)(d)=
          2\int_{|z|<1}(1-|z|^2)A_E(\sqrt{2d}\,z)\,\dd z.
\end{equation}
The inverse-map estimates bound $A_E-A_{\widetilde E}$ in each fixed
$C^s$ norm by a sufficiently high finite norm of the two input
differences. The integral in \eqref{eq:v9-fixed-disk}, as a function
of $t=\sqrt{2d}$, is smooth and even. A smooth even function $h(t)$
has a smooth right representation in $t^2$ with its first $k$
derivatives controlled by derivatives of $h$ through order $2k$.
For example $h'(t)/t=\int_0^1h''(st)\,\dd s$; iteration, applied
also to the even function $h''$, proves the assertion and its
difference estimate. Parameter derivatives commute with these
integrals. Applying it to \eqref{eq:v9-fixed-disk} proves the smooth
extension and \eqref{eq:v9-morse-transport}. All sublevel sets lie
inside the common chart by the uniform quadratic lower bound, so
there is no additional chart-boundary term.
\end{proof}

Apply this proposition to $(E_j,b_j)$ and
$(S_0\oplus S_p,B_0\otimes B_p)$. The exact physical multiplier is
$q_p/(2\pi A\sinh(j\gamma))$, by
\eqref{eq:g-full-flux}. Thus the exponentially small flux is retained
\emph{outside} a uniformly stable normalized integral. In this form
the proof yields all fixed-order derivatives on a closed collar and
the rare-event-preserving experiment transfer, rather than only an
identity at the normal orbit.

The operator identities above, the trace estimate and the smooth
integration principle are reusable elementary ingredients. The
billiard-specific work is to verify their simultaneous hypotheses for
physical stationary records with arbitrary collision order. The
resulting law has, in addition to the experiment comparison already
proved, the nonlinear parity consequence established next. This
consequence concerns the complete boundary functions at positive
offset and does not follow from a finite list of leading amplitudes.
